\documentclass[10pt,a4paper]{amsart}
\usepackage[margin=19mm]{geometry}
\usepackage{amsmath,amssymb,mathtools}
\usepackage[scr=rsfs,cal=euler]{mathalfa}
\usepackage{xfrac}
\usepackage[unicode,hidelinks,bookmarks=false]{hyperref}
\newcommand{\ie}{i.e.\ }
\newcommand{\cf}{cf.\ }
\newcommand{\resp}{respectively\ }
\newcommand{\Subsection}{\subsection}
\providecommand{\nobreakdash}{\nobreak\hbox{-}}
\setlength{\emergencystretch}{2em}
\multlinegap0pt
\usepackage{bm}
\usepackage{bbm}
\usepackage{dsfont}
\usepackage{xspace}
\usepackage{cancel}
\usepackage{mathtools}
\usepackage[shortlabels]{enumitem}
\usepackage{amsthm}
\makeatletter
\@ifundefined{lemma}{\newtheorem{lemma}{Lemma}}{}
\@ifundefined{theorem}{\newtheorem{theorem}{Theorem}}{}
\makeatother
\newcommand{\AFGPM}{\sc Anti-ForcingGPM}
\newcommand{\AFS}{\rm anti-forcing~set}
\newcommand{\FGPM}{\sc ForcingGPM} % Forcing Set for a given Perfect Matching
\newcommand{\FS}{\rm forcing~set}
\title[Hardness of Forcing Unique Perfect Matchings in Bipartite Gr]{Hardness of Forcing Unique Perfect Matchings in Bipartite Graphs of Maximum Degree 3}
\author{Ryoma Aoshima, Takashi Horiyama, Atsuki Nagao, Fumiya Sakamoto, Hibiki Sato, Kazuhisa Seto, Karin Umebayashi}
\date{Source: arXiv:2608.18617v1 (CC BY 4.0)}
\begin{document}
\maketitle

\noindent\emph{Source and licence.} Hardness of Forcing Unique Perfect Matchings in Bipartite Graphs of Maximum Degree 3. Version arXiv:2608.18617v1 (CC BY 4.0), distributed under the Creative Commons Attribution 4.0 International licence, as recorded in the arXiv licence metadata for this exact version and shown on the abstract page https://arxiv.org/abs/2608.18617v1. Full rights notice: licenses/arxiv-2608.18617-CC-BY-4.0.txt.\par\smallskip

\noindent\textit{Editorial scope.} Counted unit: the Theorem whose source label is \texttt{thm:FGPMtoAFGPM\_reduction} in the article (source0.tex lines 288--290, source label \texttt{thm:FGPMtoAFGPM\_reduction}), delivered together with the article's own complete original proof of it (source0.tex lines 292--321, one \texttt{proof} environment of 30 lines). No mathematics is written, repaired or re-derived: statement and proof are the article's own text line for line. Required context retained uncounted: lem:one-to-one (lines 202--204); lem:replace\_A (lines 260--262). Every definition, display and label of those lines is retained unchanged. Omitted from this package and not redistributed: 1--201, 205--259, 263--287, 291--291, 322--558. Nothing inside a retained excerpt is omitted or altered: every retained body is delivered exactly as the article's lines are written, so no omission receipt of any kind accompanies this package. Citations: the retained excerpts carry no citation command at all, so no bibliography entry of the article is transcribed and no reference is redistributed. Reference adaptation: none was needed and none was made; every reference inside the delivered unit resolves inside it, so no unresolved reference or citation is produced. Printed slips kept literal: no printed word, symbol, hypothesis or number of the counted statement or of its proof was altered. Layout and class: the article's own class is \texttt{article} with packages including fullpage, graphicx, fontenc, lineno, amsmath, amssymb, mathtools, amsthm, url, comment, fancybox, ascmac, subcaption, color; the wrapper is the lane's pinned \texttt{amsart} editorial wrapper at 10pt on A4, which restates the theorem-like environments the retained text needs and adds the article's own macro definitions that the retained text needs (\texttt{\textbackslash AFGPM}, \texttt{\textbackslash AFS}, \texttt{\textbackslash FGPM}, \texttt{\textbackslash FS}), with extra wrapper packages (bm, bbm, dsfont, xspace, cancel, mathtools, enumitem). The delivered numbering is the unit's own. Assets: no retained span contains a figure environment, a graphics-inclusion command, a PDF-page inclusion, a table or a TikZ picture; no image file is copied into this package, the package declares no asset dependency and no image file is redistributed with it. Licence: version arXiv:2608.18617v1 of this article is distributed under the Creative Commons Attribution 4.0 International licence, as recorded in the arXiv licence metadata for this exact version and shown on the abstract page https://arxiv.org/abs/2608.18617v1. A full rights notice is delivered at licenses/arxiv-2608.18617-CC-BY-4.0.txt. Characters: every retained excerpt is pure ASCII, so the delivered engine, XeLaTeX with its default Latin Modern text font, reports no missing character.
\begin{lemma} \label{lem:one-to-one}
There is a one-to-one correspondence between the perfect matchings of $G$ and those of $G'$.
\end{lemma}

\begin{lemma} \label{lem:replace_A}
For any perfect matching $M' \in \mathcal{M}(G')$ and any anti-forcing set $A$ of $M'$, there exists an anti-forcing set $A'$ of $M'$ such that $|A'| \le |A|$ and $A'\subseteq E_{\textnormal{mid}}$.
\end{lemma}

\begin{theorem}\label{thm:FGPMtoAFGPM_reduction}
There exists a polynomial-time reduction from {\FGPM} to {\AFGPM}.
\end{theorem}

\begin{proof}
We transform an instance $(G, M, k)$ of {\FGPM} into an instance $(G', \phi(M), k)$ of {\AFGPM}.
The graph $G'$ is constructed from $G$ by subdividing all edges of $G$, and $\phi$ is the mapping defined in Lemma~\ref{lem:one-to-one}.
Since each edge $e \in E(G)$ is replaced with two vertices and three edges, the transformation from $G$ to $G'$ runs in $O(|E(G)|)$ time.

We show that the following statements are equivalent:
\begin{enumerate}[(i)]
    \item There exists a {\FS} $F$ of $M$ in $G$ with $|F| \le k$.
    \item There exists an {\AFS} $A$ of $\phi(M)$ in $G'$ with $|A| \le k$.
\end{enumerate}

\noindent (i)$\Rightarrow$(ii):
For each edge $e \in F$, include the edge $e^{\textnormal{mid}} \in E^e_{\textnormal{mid}}$ in $A$.
We show that the resulting set $A$ is an {\AFS} of $\phi(M)$.
Suppose that $A$ is not an {\AFS} of $\phi(M)$.
Then there exists a perfect matching $M' \neq \phi(M)$ such that $A \subseteq E(G') \setminus M'$.
Since $\phi$ is bijective, we have $\phi^{-1}(M') \neq \phi^{-1}(\phi(M)) = M$ such that $F \subseteq \phi^{-1}(M')$, which contradicts the fact that $F$ is a {\FS} of $M$.
Hence, $A$ is an {\AFS} of $\phi(M)$, and $|A| = |F| \le k$.
Therefore, (i)$\Rightarrow$(ii) holds.\\

\noindent (ii)$\Rightarrow$(i):
Let $A$ be an {\AFS} that uniquely determines a perfect matching $M'$ of $G'$.
By Lemma~\ref{lem:replace_A}, there exists another {\AFS} $A' \subseteq E_{\textnormal{mid}}$ of $M'$ such that $|A'| \le |A|$.
We construct a {\FS} $F$ of the perfect matching $\phi^{-1}(M')$ of $G$ as follows: for each $e^{\textnormal{mid}} \in A'$, include $e$ in $F$.
Suppose that $F$ is not a {\FS} of $\phi^{-1}(M')$.
Then there exists a perfect matching $M'' \neq \phi^{-1}(M')$ such that $F \subseteq M''$.
Since $\phi$ is bijective, we have $\phi(M'') \neq \phi(\phi^{-1}(M')) = M'$ such that $A' \subseteq E(G') \setminus \phi(M'')$, which contradicts the fact that $A'$ is an {\AFS} of $M'$.
Hence, $F$ is a {\FS} of $\phi^{-1}(M')$, and $|F| = |A'| \le |A| \le k$.
Therefore, (ii)$\Rightarrow$(i) holds.
\end{proof}

\end{document}
